\begin{abstract}
When UTXO recipients wait for public execution before spending, every
payment hop adds a confirmation delay. Next lets a recipient continue after
one certification round. Our insight is that the certificate traveling with
an output names who answers for a missing source: its direct issuer.
Members of the consuming organization lock inputs and reserve capacity for
all certified outputs; a recipient that verifies their quorum certificate
(TXCer) may request a successor, subject to successor admission and
capacity. If the successor executes first, public execution consumes the
certified output and records a funded obligation against that issuer.
Source execution fulfills the obligation; otherwise an included, valid post-deadline
public decision compensates the gap from the issuer's reserve, a later
source repays it, and a source that never executes leaves the loss with the
issuer. Authorized historical repair records that debit in the consuming
input under the original block identity, without re-executing payments or making economic closure wait for
historical repair. For four-member organizations and a four-validator
committee, each with at most one static Byzantine member, we prove unique
consumption, coverage of every certificate's outstanding liability, CAL
conservation, and order-independent terminal CAL holdings for finite,
acyclic, source-closed, conflict-free payment sets whose payments each
succeed exactly once. On one host with synchronous member commits and NoSync wallets, a
100-hop chain reaches its last certified receipt in a median 4.707\,s,
versus 64.479\,s with per-hop waiting.
\end{abstract}

\begin{IEEEkeywords}
Blockchain, certified payments, collateral, Byzantine fault tolerance,
chameleon hash.
\end{IEEEkeywords}
